% =============================================================
% Matriz de Rastreabilidade de Testes — Sprint 1
% GreenHerb API  |  Engenharia de Software II
% Compilar com: pdflatex matriz_rastreabilidade.tex  (2×)
% =============================================================
\documentclass[a4paper,11pt]{article}

\usepackage[utf8]{inputenc}
\usepackage[T1]{fontenc}
\usepackage[portuguese]{babel}
\usepackage[a4paper, top=2.5cm, bottom=2.5cm, left=2.5cm, right=2.5cm]{geometry}
\usepackage{booktabs}
\usepackage{longtable}
\usepackage{array}
\usepackage[table,dvipsnames]{xcolor}
\usepackage{pdflscape}
\usepackage{hyperref}
\usepackage{parskip}
\usepackage{titlesec}
\usepackage{fancyhdr}
\usepackage{microtype}

% ------------------------------------------------------------------
% Cores
% ------------------------------------------------------------------
\definecolor{headerblue}{RGB}{68,114,196}
\definecolor{rowalt}{RGB}{235,241,250}
\definecolor{grouprow}{RGB}{217,225,242}
\definecolor{titlecolor}{RGB}{31,73,125}
\definecolor{invalida}{RGB}{255,235,238}
\definecolor{valida}{RGB}{232,245,233}

% ------------------------------------------------------------------
% Secções
% ------------------------------------------------------------------
\titleformat{\section}{\large\bfseries\color{titlecolor}}{}{0em}{}%
  [\color{headerblue}\titlerule]
\titleformat{\subsection}{\normalsize\bfseries\color{titlecolor}}{}{0em}{}

% ------------------------------------------------------------------
% Tipos de coluna (já com \footnotesize e quebra de linha automática)
% ------------------------------------------------------------------
\newcolumntype{L}[1]{>{\raggedright\arraybackslash\footnotesize}p{#1}}
\newcolumntype{C}[1]{>{\centering\arraybackslash\footnotesize}p{#1}}

\setlength{\tabcolsep}{4pt}
\renewcommand{\arraystretch}{1.4}

% ------------------------------------------------------------------
% Cabeçalho / rodapé
% ------------------------------------------------------------------
\pagestyle{fancy}
\fancyhf{}
\rhead{\footnotesize GreenHerb API — Matriz de Rastreabilidade}
\lhead{\footnotesize Sprint 1 — Módulo de Autenticação}
\cfoot{\footnotesize \thepage}
\renewcommand{\headrulewidth}{0.4pt}

\hypersetup{
  colorlinks=true,
  linkcolor=headerblue,
  urlcolor=headerblue,
  pdftitle={Matriz de Rastreabilidade — Sprint 1},
}

% =============================================================
\begin{document}
% =============================================================

% ------------------------------------------------------------------
% Título
% ------------------------------------------------------------------
\begin{center}
  {\LARGE\bfseries\color{titlecolor} GreenHerb API}\\[6pt]
  {\large\bfseries Matriz de Rastreabilidade de Testes}\\[4pt]
  {\normalsize Sprint 1 — Módulo de Autenticação}\\[12pt]
  \begin{tabular}{@{}ll@{}}
    \textbf{UC:}   & Engenharia de Software II — Testes de Software \\[3pt]
    \textbf{Data:} & 2026-05-07 \\
  \end{tabular}
\end{center}

\vspace{0.4cm}
\tableofcontents
\newpage

% =============================================================
\section{Legenda de Prefixos}
% =============================================================

\begin{tabular}{@{}C{2cm}L{9cm}@{}}
  \toprule
  \textbf{Prefixo} & \textbf{Nível de Teste} \\
  \midrule
  \rowcolor{rowalt} TU & Teste de Unidade \\
                    TI & Teste de Integração \\
  \rowcolor{rowalt} TS & Teste de Sistema \\
  \bottomrule
\end{tabular}

% =============================================================
\section{Requisito Funcional de Referência}
% =============================================================

\begin{tabular}{@{}C{1.4cm}L{13cm}@{}}
  \toprule
  \textbf{ID} & \textbf{Descrição} \\
  \midrule
  \rowcolor{rowalt}
  RF-01 & O sistema deve autenticar utilizadores com \texttt{username} e
          \texttt{password} válidos, devolvendo um \textit{access token} e um
          \textit{refresh token} JWT. Credenciais ausentes ou incorrectas devem
          ser rejeitadas. \\
  \bottomrule
\end{tabular}

% =============================================================
\section{Particionamento de Equivalência}
% =============================================================

Esta secção define as classes de equivalência identificadas para os dois
parâmetros de entrada da função de autenticação
(\texttt{login(username, password, findUser)}).

\subsection{Partições para \texttt{username}}

\begin{tabular}{@{}C{1.4cm}C{2.2cm}L{5cm}L{4.5cm}@{}}
  \toprule
  \textbf{ID} & \textbf{Classe} & \textbf{Condição} & \textbf{Representante(s)} \\
  \midrule
  \rowcolor{invalida}
  P-U1 & Inválida & Valor ausente: \texttt{null}, \texttt{undefined} ou string vazia
       & \texttt{null} ; \texttt{""} \\
  \rowcolor{invalida}
  P-U2 & Inválida & String não vazia mas sem utilizador correspondente na BD
       & \texttt{"desconhecido"} \\
  \rowcolor{valida}
  P-U3 & Válida   & Utilizador existente na BD
       & \texttt{"alice"} \\
  \bottomrule
\end{tabular}

\subsection{Partições para \texttt{password}}

\begin{tabular}{@{}C{1.4cm}C{2.2cm}L{5cm}L{4.5cm}@{}}
  \toprule
  \textbf{ID} & \textbf{Classe} & \textbf{Condição} & \textbf{Representante(s)} \\
  \midrule
  \rowcolor{invalida}
  P-P1 & Inválida & Valor ausente: \texttt{null}, \texttt{undefined} ou string vazia
       & \texttt{null} ; \texttt{""} \\
  \rowcolor{invalida}
  P-P2 & Inválida & String fornecida que não corresponde ao \textit{hash} armazenado
       & \texttt{"errada"} \\
  \rowcolor{valida}
  P-P3 & Válida   & Password correcta para o utilizador
       & \texttt{"senha123"} \\
  \bottomrule
\end{tabular}

% =============================================================
\section{Matriz de Rastreabilidade}
% =============================================================

% Larguras (landscape A4 com margens 2,5 cm ≈ 24,7 cm disponíveis):
%   0,9 + 1,1 + 1,0 + 1,0 + 3,0 + 4,5 + 5,5 + 5,5 = 22,5 cm
%   separadores: 7 × 8pt ≈ 2,0 cm  →  total ≈ 24,5 cm  ✓

\begin{landscape}
\setlength{\tabcolsep}{4pt}
\renewcommand{\arraystretch}{1.5}

\begin{longtable}{
  C{0.9cm}   % ID CT
  C{1.1cm}   % Requisito
  C{1.0cm}   % Part. username
  C{1.0cm}   % Part. password
  L{3.0cm}   % Endpoint
  L{4.5cm}   % Técnica
  L{5.5cm}   % Resultado Esperado
  L{5.5cm}   % Pré-condições
}

\caption{Casos de Teste — Sprint~1, função \texttt{login} (Particionamento de Equivalência)}
\label{tab:matriz}\\

% --- Cabeçalho primeira página ---
\toprule
\textbf{\footnotesize ID CT} &
\textbf{\footnotesize Req.} &
\textbf{\footnotesize Part. \textit{user}} &
\textbf{\footnotesize Part. \textit{pass}} &
\textbf{\footnotesize Endpoint Exercitado} &
\textbf{\footnotesize Técnica Aplicada} &
\textbf{\footnotesize Resultado Esperado} &
\textbf{\footnotesize Pré-condições} \\
\midrule
\endfirsthead

% --- Cabeçalho páginas seguintes ---
\multicolumn{8}{l}{\footnotesize\textit{(continuação da Tabela~\ref{tab:matriz})}}\\[3pt]
\toprule
\textbf{\footnotesize ID CT} &
\textbf{\footnotesize Req.} &
\textbf{\footnotesize Part. \textit{user}} &
\textbf{\footnotesize Part. \textit{pass}} &
\textbf{\footnotesize Endpoint Exercitado} &
\textbf{\footnotesize Técnica Aplicada} &
\textbf{\footnotesize Resultado Esperado} &
\textbf{\footnotesize Pré-condições} \\
\midrule
\endhead

\midrule
\multicolumn{8}{r}{\footnotesize\textit{continua na página seguinte}}\\
\endfoot

\bottomrule
\endlastfoot

% ------------------------------------------------------------------
% TU-01
% ------------------------------------------------------------------
\rowcolor{invalida}
TU-01 & RF-01 & P-U1 & P-P3 &
\texttt{POST /auth/login} (unidade, indirecto) &
PE — classe inválida de \texttt{username} (\texttt{null}) &
Lança \texttt{Error}: \textit{"Username and password are required"} &
\texttt{JWT\_SECRET} definido; \texttt{findUser} não deve ser invocado \\

% ------------------------------------------------------------------
% TU-02
% ------------------------------------------------------------------
\rowcolor{invalida}
TU-02 & RF-01 & P-U1 & P-P3 &
\texttt{POST /auth/login} (unidade, indirecto) &
PE — classe inválida de \texttt{username} (string vazia \texttt{""}) &
Lança \texttt{Error}: \textit{"Username and password are required"} &
Idem \\

% ------------------------------------------------------------------
% TU-03
% ------------------------------------------------------------------
\rowcolor{invalida}
TU-03 & RF-01 & P-U2 & P-P3 &
\texttt{POST /auth/login} (unidade, indirecto) &
PE — classe inválida de \texttt{username} (utilizador inexistente) &
Lança \texttt{Error}: \textit{"Invalid credentials"} &
\texttt{findUser} mockado a devolver \texttt{null} para o \textit{username} fornecido \\

% ------------------------------------------------------------------
% TU-04
% ------------------------------------------------------------------
\rowcolor{invalida}
TU-04 & RF-01 & P-U3 & P-P1 &
\texttt{POST /auth/login} (unidade, indirecto) &
PE — classe inválida de \texttt{password} (\texttt{null}) &
Lança \texttt{Error}: \textit{"Username and password are required"} &
\texttt{JWT\_SECRET} definido; \texttt{findUser} não deve ser invocado \\

% ------------------------------------------------------------------
% TU-05
% ------------------------------------------------------------------
\rowcolor{invalida}
TU-05 & RF-01 & P-U3 & P-P1 &
\texttt{POST /auth/login} (unidade, indirecto) &
PE — classe inválida de \texttt{password} (string vazia \texttt{""}) &
Lança \texttt{Error}: \textit{"Username and password are required"} &
Idem \\

% ------------------------------------------------------------------
% TU-06
% ------------------------------------------------------------------
\rowcolor{invalida}
TU-06 & RF-01 & P-U3 & P-P2 &
\texttt{POST /auth/login} (unidade, indirecto) &
PE — classe inválida de \texttt{password} (incorrecta) &
Lança \texttt{Error}: \textit{"Invalid credentials"} &
\texttt{findUser} mockado com utilizador válido; \textit{hash} de \textit{bcryptjs} gerado em \texttt{beforeAll} (salt\,=\,10) \\

% ------------------------------------------------------------------
% TU-07
% ------------------------------------------------------------------
\rowcolor{valida}
TU-07 & RF-01 & P-U3 & P-P3 &
\texttt{POST /auth/login} (unidade, indirecto) &
PE — classe válida (\texttt{username} e \texttt{password} correctos) &
Retorna objecto \texttt{\{accessToken, refreshToken, perfil\}}; tokens verificáveis com \texttt{JWT\_SECRET} &
\texttt{findUser} mockado com \texttt{\{\_id, perfil, password\}} correctos; credenciais coincidem \\

\end{longtable}
\end{landscape}

% =============================================================
\section{Resultados de Execução}
% =============================================================

Comando: \texttt{npx jest tests/unit/authService.test.js {-}{-}verbose}

\vspace{0.4cm}

\begin{tabular}{@{}C{1.2cm}L{9.5cm}C{1.8cm}@{}}
  \toprule
  \textbf{ID CT} & \textbf{Descrição do teste} & \textbf{Estado} \\
  \midrule
  \rowcolor{invalida}
  TU-01 & \textit{username} null (P-U1): lança erro de campo obrigatório        & PASS \\
  \rowcolor{invalida}
  TU-02 & \textit{username} vazio (P-U1): lança erro de campo obrigatório       & PASS \\
  \rowcolor{invalida}
  TU-03 & \textit{username} inexistente (P-U2): lança erro de credenciais inválidas & PASS \\
  \rowcolor{invalida}
  TU-04 & \textit{password} null (P-P1): lança erro de campo obrigatório        & PASS \\
  \rowcolor{invalida}
  TU-05 & \textit{password} vazia (P-P1): lança erro de campo obrigatório       & PASS \\
  \rowcolor{invalida}
  TU-06 & \textit{password} incorrecta (P-P2): lança erro de credenciais inválidas & PASS \\
  \rowcolor{valida}
  TU-07 & Credenciais válidas (P-U3 $\times$ P-P3): retorna tokens e perfil     & PASS \\
  \midrule
  \multicolumn{2}{@{}l}{\textbf{Total}} & \textbf{7 / 7} \\
  \bottomrule
\end{tabular}

\end{document}
